% =====================================================
\section{Conclusion}
\label{sec:conclusion}
% =====================================================

We compared classical and learned regularization for sparse-view CT, using a noise-conditioned DnCNN in Plug-and-Play proximal gradient against FBP, Tikhonov and TV, under noise, undersampling, forward-model mismatch, distribution shift and a lesion-detection task, first on $28\times28$ digits with Gaussian and Poisson noise and then on real LoDoPaB-CT slices.

On digits the learned prior is about $3.4$\,dB better than box-constrained TV at moderate noise ($+3.8$\,dB over TV-ADMM). The advantage shrinks at both ends of the noise range: at $\sigma=0$ TV is ahead under every data model tested, and at $\sigma=0.01$ the ranking depends on the data model and the TV solver. The advantage halves for a near shift (EMNIST letters) and reverses for a far shift (Fashion-MNIST, Shepp--Logan phantom). The matched-prior control, in which only the denoiser weights change, shows that the reversal is a property of the prior and not of the PnP scheme, and the error decomposition shows that most of the excess error lies in the well-determined directions, not only in the null space.
PSNR is not task performance: in distribution the PSNR-tuned prior attenuates small lesions and detects them less reliably than TV, and under mismatch it creates more false lesions than box-constrained TV for seven of eight lesion types. Re-tuning trades these against each other, and an operating point competitive with TV on PSNR, detection and false lesions exists only for the Fashion-trained prior, which we hypothesise reflects the variety of its training data (untested with two priors).
The number of iterations acts as a regularisation parameter: the data residual stays at the noise level while the error grows in weakly determined directions, which we interpret as semi-convergence but did not test.

On real CT slices the PnP prior does not consistently improve on a well-tuned TV run for 300 iterations in PSNR (differences between $-0.93$ and $+0.33$\,dB over nine dose and angle conditions) and has a lower SSIM in eight of nine; the two are not matched in cost or in tuned parameters. It detects small high-contrast lesions better than that TV in all nine conditions (significant in six, negligible against the 150-iteration TV at full dose for $K=30$ and $60$) but not better than unregularised FBP, whereas the PSNR-tuned TV detects them worse than FBP in all nine. Its false-blob counts on real slices are below those of FBP-Hann in the four conditions tabulated and not systematically different from those of TV-long. Under a shift to non-anatomical phantoms its PSNR and SSIM deficit against TV-long is consistent, its detection advantage is significant in one of six conditions, and its false-blob count relative to TV changes sign with the number of angles.

Our conclusions concern one small denoiser architecture, two training distributions on digits and one CT training run. They are conditional on the regime, the data model, the TV variant, the tuning protocol and a location-known detection statistic at PSNR-tuned operating points, and the CT analysis was developed after seeing results, without a pre-specified primary comparison.
They suggest evaluating learned priors against baselines with the same constraints and a tuned or compute-matched iteration count, with hyperparameters tuned on held-out data, several training seeds, mismatched and matched priors, and task-based metrics next to PSNR and SSIM.
Future work should designate the primary analyses in advance and use all 60 test patients with several noise realisations; add paired intervals for the digit detection differences; check the PnP--TV difference on real data and not only the TV baseline; tune TV and PnP for detection on CT; use a finer TV and Hann grid, a compute-matched TV, several denoiser seeds, stronger denoisers and the full training set; run box-constrained TV in every experiment; and use clinically calibrated lesion contrasts with model observers, re-tuning under forward-model mismatch, an iteration curve on noise-free data, and provably convergent PnP variants.